\documentclass[12pt]{article}

\usepackage[a4paper, total={6in, 8in}]{geometry}
\usepackage{textcomp}

\begin{document}
\setlength{\parindent}{0pt}

\def \t {\textrightarrow}
\def \v {\vspace{1em}}
\def \bi {\begin{itemize}}
\def \ei {\end{itemize}}
\def \s[#1] {\section*{#1}}
\def \ss[#1] {\subsection*{#1}}
\def \sss[#1] {\subsubsection*{#1}}

\s[Montale]
\ss[Spesso il male di vivere ho incontrato]
Negli ossi di seppia \\
Il male di vivere è il volto del suo dramma interiore \t il dramma della mancanza di liberta \\
Il testo immette nella definizione degli oggetti del male di vivere \t l'oggettistica montaliana \\
Il male di vivere è incarnata negli oggetti, ma non sappiano da dove derivi questo smarrimento dell'auotre \\
Onomatopea al verso 2 \\
Strozzato è personificazione non proprio \t anche sinestesia, passabile metafora \t qua la sinestesia incarna il correlativo oggettivo \\
Anafora di era \\
Le foglie le abbiamo viste in soffia il vento e nevica la frasca \\
Presente la lettera "r", suono che graffia \t anaofra ripetuta e correttiva (prima mette riarsa e poi era il cavallo) con l'anastrofe \\
L'essere stramazzato percepisce la tensione \t si legge in questo la sfinitezza di questo cavallo \\
La prima parte descrive l imperfetto \t azione non conclusa \t quindi il male di vivere è ancora presente, non appartiene al passato \\
Divina (miniscola) indifferenza (maiuscola) \t dio è indifferente, non pensa a noi \t qua l'indifferenza si concrettiza in una statua assopita, mentre la nuvola è lontana e scollata dall'uomo \\
C'è un qualcosa di lucreziano in questa lirica \t distacco di dio, gli dei sono lontani e indifferenti \t come il saggio guarda dall'alto della sua torre d'avorio \\
Non sappiamo l'emozione che l ha portato a usare questi oggetti pero \\
Il vivere è il male di vivere, ottica esistenzialista è presente

\v

Importanza del ricordo, che lo corella a leopardi \t Clizia stessa rappresenta l'illuminazione e la svolta, ma il distacco da lei lo segnera per tutta la vita \\
Lo sai che debbo riperderti e non posso \t lei è come l'alba per un nuovo giorno

\ss[Cigola la carrucola del pozzo]
Casa rurale \t il pozzo serve per somministrare acqua \\
Narciso, lo specchiarsi nell'acqua \\
Acqua dal fondo risale alla luce, e si fonde con la luce \\
Un ricordo trema: sinestesia \t l'immagine che appare scatena il ricordo \t trema perchè è riflesso nell'acqua che si muove \\
Come le labbra, come il volto \t l'immagine di noi, il passato è stravolto \t noi non siamo quello che eravamo \\
Stride \t predilezione per le onomatopea (gia con cigola) \t immagine è stridente, evoca qualcosa di arrugginito, poco fluido \t un movimento a fatica, come il ricordo che a fatica \\
Il ricordo non è consolatorio, come in leopardi \t + vicino al leopardi di "A se stesso", in cui il ricordo non è consolatorio \\
Lo stravolgimento del passato \t qua abbiamo un altra persona \\
La ruota ti all' \t colloquilatia \\
Qualcosa di rimato è rimasto \t non c'è abbandono delle rime \t una distanza ci divide, immagine di desolazione e sofferenza \t come una piaga aparta dell'anima \t dolore che non trova consolazione \\
Come gli ossi di seppia, che sono cartilagini facilmente scheggiabili \\
Il tema del ricordo non è il solo \t desolazione, abbandono, vuoto, alienazione da se stessi \t emozione che si fa oggetto non ci è svelata \\

\ss[Non recidere, forbice, quel volto]
Volto della memoria, con un non so che di evocativo qua \\
Evocazione iniziale quasi pascoliana \t che ricordano anche il d'annunzio + intimistico dei madrigali d'estate, e un pascoli di Novembre \\
Le occasioni del 39, parola chiave Clizia \\
Montala capisce che nel male ci sono delle occasioni \\
Liirica lapidaria, sintetica, rimate secondo una certa tradizione \t ma la negazione porta indietro a non chiedere la parola \\
Sfollare = liberarsi \\
Non verso 1 e 3 \t anafora \\
La nebbia della mente \t ognuno sa quanto sia difficile configurare i tratti di qualcuno che non c'è piu \\
Il testo lascia una consapevolezza, il non ricordare una persona ma di ricordare cio che continua a vivere dentro di me di essa \\
Un colpo dura \t è sinestesia \t mentre svettare è montalismo, qua vuol dire tranciare la cima di un'acacia \t colpo duro trancia, come una vita tranciata \\
L'acacia evoca anche il miele \t acacia ferita, non tanto con una personificazione ma con una sinestesia \t cade \\
Acacia fa cadere, scrolla il guscio di cicala \\
Novembre evoca pascoli, belletta d'annunziano ma usato in un altro contesto \t che parte con reminescenze, attualizza il tema del ricordo \\
La nebbia guaagna terreno \t mentre pascoli parte da un lutto
\end{document}
